\section{Conclusion}

We present a self-improving framework for training vision-language model judges that eliminates the need for human preference annotations through strategic synthetic data generation and iterative refinement. Our approach demonstrates that a compact 11B parameter model can achieve competitive performance with significantly larger proprietary systems across multiple evaluation dimensions, particularly excelling in general instruction following, hallucination detection, and VQA tasks. Our analysis reveals clear pathways for further improvement through targeted data augmentation in specialized domains such as safety evaluation, while our core methodology remains applicable to emerging visual tasks where ground-truth annotations are scarce or unavailable.

\section*{Limitations}

\xhdr{Safety Evaluation Limitations} The modest improvements observed on safety benchmarks reflect a deliberate methodological boundary: our framework does not generate synthetic training data containing biased, toxic, or otherwise harmful content. Robust safety evaluation requires specialized infrastructure including red-teaming protocols, carefully designed adversarial examples, and nuanced understanding of subtle policy violations--components that fall outside the scope of our general-purpose self-improvement approach. Developing effective safety evaluation capabilities within self-improving frameworks remains an important open challenge, one that will require principled approaches for safely generating and learning from examples of harmful content without amplifying risks.

\xhdr{Task-Specific Performance Variation} Our synthetic data generation framework exhibits differential effectiveness across evaluation dimensions, reflecting the inherent complexity and diversity of multimodal judgment tasks. While substantial improvements emerge in general instruction following, hallucination detection, and visual question answering, other capabilities such as complex reasoning and nuanced preference judgments show more limited gains. This performance variance suggests potential benefits from mixture-of-experts architectures, where a routing mechanism could direct different judgment queries to specialized expert judges trained on domain-specific synthetic data. Exploring such multimodel approaches represents a promising direction for future work that could build upon our iterative self-training framework, extending the principle of self-generated supervision signals to specialized judge models tailored for different evaluation dimensions